\section{Convergent generating operators for the selected radial diagonals}
\label{sec:ns-diagonal-generator}

This section records the exact generating operator for the highest profile-order
diagonals of the source radial jets.  It is a selected diagonal, not the full
Navier--Stokes field.  The axis data, viscosity, coordinates, cutoffs and
source parameters are retained exactly as in Sections~\ref{sec:ns-actual-witness}
and~\ref{sec:ns-finite-radial-recursion}; the source existence theorem remains
an imported result~\cite{openaiNS2026release}.  The scalar kernel below is the
source's Appendix~B function (B.11), and its Bessel provenance is recorded in
the NIST DLMF~\cite[\S10.2,\S10.9]{olverDLMF2}.

\subsection{Local analytic axis data and their radius}

Use the source parameters
\[
 A=\tfrac12+h,\quad D=\tfrac12-h,\quad k=1+h,\quad 0<h<\tfrac1{100},
 \quad \nu>0,
\]
and put $\tau=1-t$, $\sigma=r^2/2$, $y=z/\sqrt\nu$,
\[
  \tau=q(1-\eta^2),\quad\quad y=q^D\eta.
\]
On the central preterminal slab, where the source cutoffs are one near the
axis, the retained axis data are
\begin{equation}\label{eq:ns-diagonal-axis-data}
 w_0(z,t)=\sqrt\nu\,q^{-A}U(\eta),\qquad
 b_1(z,t)=\frac{2}{C}q^{-k}\phi(\eta).
\end{equation}
Here $U$ and the holomorphic germ $\phi$ are the source functions defined in
Appendix~B; no new amplitude or normalization is introduced.

For a real preterminal point $(z_0,t_0)$, write
\[
 F(\eta)=\eta(1-\eta^2)^{-D},\qquad
 z=\sqrt\nu\,\tau^D F(\eta),\qquad q=\frac{\tau}{1-\eta^2}.
\]
If $\eta_0$ is the corresponding real parameter, then
\[
 F'(\eta_0)=(1-\eta_0^2)^{-D-1}
 \bigl(1-2h\eta_0^2\bigr)>0.
\]
Choose a disk avoiding $\pm1$, the roots of $1-2h\eta^2$, and the poles of
the rational logarithmic derivative defining $\phi$.  On that disk the
positive real branch of $(1-\eta^2)^{-D}$ has a holomorphic logarithm.  The
map $F$ has a quantitative inverse there: if $a=F'(\eta_0)$ and $M_2$ is
any bound for $|F''|$ on the disk, then
\[
  \delta\leq \frac{a}{4(M_2+1)},\quad\quad
 R=\frac{\sqrt\nu\,\tau_0^D a\delta}{4}
\]
give a disk $|z-z_0|<R$ on which the contraction
\[
 T_Y(\eta)=\eta+\frac{Y-F(\eta)}a,\quad\quad
 Y=\frac{z}{\sqrt\nu\,\tau_0^D},
\]
has a unique holomorphic fixed point.  Thus both functions in
\eqref{eq:ns-diagonal-axis-data} have a positive local holomorphic radius,
with the factor $\sqrt\nu\,\tau_0^D$ retained.  No uniform radius up to
$t=1$ is asserted.

\subsection{The two generating kernels}

Let $f$ be holomorphic on $|\zeta-z_0|<R$.  Whenever
$|z-z_0|+|r|<R$, define
\begin{align}
 H_0[f](r,z)&=\frac1{2\pi}\int_0^{2\pi}
 f(z+ir\cos\theta)\,d\theta,\label{eq:nsdiag-H0}\\
 H_1[f](r,z)&=\frac1\pi\int_0^{2\pi}
 \sin^2\theta\,f(z+ir\cos\theta)\,d\theta.\label{eq:nsdiag-H1}
\end{align}
The second normalization is $1/\pi$ for the interval $[0,2\pi]$ (equivalently
$2/\pi$ on $[0,\pi]$).  Both maps are even in $r$ and hence are holomorphic
in $\sigma=r^2/2$ near the axis.  Taylor expansion and
\[
 \frac1{2\pi}\int_0^{2\pi}\cos^{2j}\theta\,d\theta
 =\frac{(2j)!}{4^j(j!)^2},\quad\quad
 \frac1\pi\int_0^{2\pi}\sin^2\theta\cos^{2j}\theta\,d\theta
 =\frac{(2j)!}{4^j j!(j+1)!}
\]
give the absolutely convergent series
\begin{align}
 H_0[f]&=\sum_{j\geq0}\alpha_j[f]\sigma^j,&
 \alpha_j[f]&=\frac{(-1)^j f^{(2j)}}{2^j(j!)^2},\label{eq:nsdiag-H0-series}\\
 H_1[f]&=\sum_{j\geq0}\beta_j[f]\sigma^j,&
 \beta_j[f]&=\frac{(-1)^j f^{(2j)}}{2^j j!(j+1)!}.
 \label{eq:nsdiag-H1-series}
\end{align}
Indeed, Cauchy's estimate $|f^{(2j)}(z)|\leq(2j)!M/a^{2j}$ and
${2j\choose j}\leq4^j$ bound the $j$th terms by a geometric series in
$|r|/a$.  The same argument on a slightly larger compact disk justifies
all displayed derivatives.  In Bessel notation the identities are
\[
 H_0=J_0(r\partial_z),\qquad
 H_1=\frac{2J_1(r\partial_z)}{r\partial_z}
 =f_0(\sigma\partial_z^2),\quad\quad
 f_0(s)=\sum_{j\geq0}\frac{(-s/2)^j}{j!(j+1)!},
\]
where the quotient is only notation for the convergent entire series; no
inverse of $\partial_z$ is used.

\subsection{PDEs, inverse maps, and the selected vector field}

Set
\[
 L_0=2\sigma\partial_{\sigma\sigma}+2\partial_\sigma+\partial_{zz},
 \qquad
 L_1=2\sigma\partial_{\sigma\sigma}+4\partial_\sigma+\partial_{zz}.
\]
The coefficient recurrences in \eqref{eq:nsdiag-H0-series}--\eqref{eq:nsdiag-H1-series}
prove, on the common analytic domain,
\begin{equation}\label{eq:ns-H-pdes}
 L_0H_0[f]=0,\qquad L_1H_1[f]=0,\quad\quad
 \partial_\sigma(\sigma H_1[f])=H_0[f],\quad\quad
 \partial_\sigma H_0[f]=-\tfrac12H_1[f''].
\end{equation}
Restriction to $\sigma=0$ is the inverse of each $H_i$ on its holomorphic
solution-germ range.  Thus these are explicit linear bijections of analytic
germs, not a claim about arbitrary smooth Taylor data.

Define
\begin{align}
 P[f]&=\left(-\frac{x_1}{2}H_1[f_z],-\frac{x_2}{2}H_1[f_z],H_0[f]\right),\\
 S[g]&=\left(-\frac{x_2}{2}H_1[g],\frac{x_1}{2}H_1[g],0\right),
 \qquad \sigma=\frac{x_1^2+x_2^2}{2}.
\end{align}
Then direct differentiation using \eqref{eq:ns-H-pdes} gives
\begin{equation}\label{eq:ns-selected-vector-identities}
 \operatorname{div}P[f]=\operatorname{div}S[g]=0,\quad\quad
 \operatorname{curl}P[f]=0,\quad\quad
 \operatorname{curl}S[g]=P[g].
\end{equation}
Every Cartesian component is harmonic because
\[
 \Delta H_0[f]=L_0H_0[f],\qquad
 \Delta(x_iG)=x_iL_1G\quad(i=1,2).
\]
The selected diagonal of the source is therefore
\begin{equation}\label{eq:ns-selected-diagonal}
 u_{\mathrm{diag}}=P[w_0]+S[b_1],
\end{equation}
with
\[
 u_{z,\mathrm{diag}}=H_0[w_0],\qquad
 u_{r,\mathrm{diag}}=-\frac r2H_1[(w_0)_z],\qquad
 u_{\theta,\mathrm{diag}}=\frac r2H_1[b_1].
\]
The inverse axis map is explicit:
\[
 w_0=u_z|_{r=0},\qquad b_1=(\operatorname{curl}u)_z|_{r=0}
 =2\lim_{r\to0}\frac{u_\theta}{r}.
\]
This is a bijection only in the stated axisymmetric, axis-regular,
componentwise-harmonic analytic-germ category.  It is not an identification
of the complete source field with a harmonic or irrotational field.

\subsection{Restoring the source cutoffs}

Let $\rho_0=1$ and $\rho_j(z,t)=\chi(c_jq(z,t))$ for $j\geq1$, exactly as
in the source.  The correct selected cutoff operators are
\[
 W_\rho[f]=\sum_{j\geq0}\rho_j\alpha_j[f]\sigma^j,
 \qquad G_\rho[f]=\sum_{j\geq0}\rho_j\beta_j[f]\sigma^j,
 \qquad B_\rho=\sigma G_\rho[b_1].
\]
The sum is locally finite for every preterminal point.  The radial component
must be reconstructed from incompressibility:
\begin{equation}\label{eq:ns-cutoff-radial}
 R_\rho[f]=r\,u_{r,\rho}[f]=-\int_0^\sigma\partial_zW_\rho[f](s,z)\,ds
 =-\sum_{j\geq0}\frac{\partial_z(\rho_j\alpha_j[f])}{j+1}\sigma^{j+1}.
\end{equation}
Dropping $\partial_z\rho_j$ would destroy divergence-freeness.  For $a=0,1$,
Let $d_{0,j}=\alpha_j$, $d_{1,j}=\beta_j$ and $H_{a,\rho}$ for the
corresponding cutoff sum.  The exact residual is
\begin{equation}\label{eq:ns-cutoff-residual}
 L_aH_{a,\rho}[f]=\sum_{j\geq0}\sigma^j\left[
 (\rho_j-\rho_{j+1})\partial_z^2d_{a,j}[f]
 +2(\partial_z\rho_j)\partial_zd_{a,j}[f]
 +(\partial_{zz}\rho_j)d_{a,j}[f]\right].
\end{equation}
The adjacent-index term and both cutoff derivatives are part of the exact
formula.  For every fixed $J$, the first $J+1$ coefficients agree with the
uncut kernels on the open region where $\rho_0,\ldots,\rho_J=1$; no region
uniform in $J$ is asserted.  Equation \eqref{eq:ns-cutoff-radial} together
with \eqref{eq:ns-cutoff-residual} gives the complete selected cutoff field,
not the full lower-profile-order source array.

\subsection{Typed handoff to the Mellin hierarchy}

Here is a compact representative of the actual local calculation. On a
compact preterminal parameter patch choose $\epsilon>0$ inside the
common radial neighborhood, and let $\kappa_\epsilon(\sigma)$ be smooth,
one for $0\leq\sigma\leq\epsilon$ and zero for $\sigma\geq2\epsilon$.
An explicit such function and the full extension of $\mathcal A_h$ to
nonzero axis values are proved in Section~\ref{sec:ns-all-axis-mellin}.
Set
\[
 a_W=\kappa_\epsilon W_\rho[w_0],\qquad
 a_G=\kappa_\epsilon G_\rho[b_1],\qquad
 a_B=\sigma a_G.
\]
The first two inputs are compact and smooth but generally do not vanish
at the axis; $a_B\in\mathcal D_0$. The original correction cutoffs
$\rho_j$ are unchanged. The coefficient maps are exactly
\begin{equation}\label{eq:nsdiag-mellin-coefficients}
 \begin{split}
 [\sigma^j]a_W&=\rho_j\alpha_j[w_0],\qquad
 [\sigma^j]a_G=\rho_j\beta_j[b_1],\quad j\geq0,\\
 [\sigma^n]a_B&=\rho_{n-1}\beta_{n-1}[b_1],\quad n\geq1.
 \end{split}
\end{equation}
The Mellin map acts on the arithmetic transform in the physical variable;
$W_h(s)$ multiplies its Mellin transform, not the physical input:
\[
 \mathfrak G_a(s)=\mathcal M(\mathcal A_h a)(s)
                =W_h(s)\mathcal M a(s).
\]
For $n\geq1$ retain
\[
 c_{h,n}=
 \begin{cases}
 W_h(-n),&n\ {\rm odd},\\
 W_h'(-n),&n\ {\rm even},
 \end{cases}
 \qquad
 E_n(\mathfrak G_a)=
 \begin{cases}
 {\rm Res}_{s=-n}\mathfrak G_a(s),&n\ {\rm odd},\\
 \mathfrak G_a(-n),&n\ {\rm even},
 \end{cases}
 \]
The proved continuation gives
\begin{equation}\label{eq:nsdiag-mellin-inverse}
 E_n(\mathfrak G_a)=c_{h,n}a_n,\qquad
 a_n=\frac{E_n(\mathfrak G_a)}{c_{h,n}}.
\end{equation}
For the two nonvanishing inputs also put
$E_0=\operatorname{Res}_{s=0}$ and $c_{h,0}=W_h(0)=(2^h-1)/4$;
then $E_0(\mathfrak G_a)=c_{h,0}a(0)$. The next section proves this
extension, its full inverse, and independence of the auxiliary localization
at the level of all residue--value coordinates. The actual angular field
uses the shifted index in the second line of
\eqref{eq:nsdiag-mellin-coefficients}, not the unshifted $G_\rho$ index.
The full source remainder is carried through the same transform below.
A selected coefficient, a regular value at a trivial zero, or a residue
is not by itself a value of the complete Weil quadratic form.

\subsection{Arithmetic interface and nonclaims}

The selected analytic field can be composed with the already proved physical
Mellin map in Sections~\ref{sec:ns-source-heat-arithmetic}--\ref{sec:ns-axis-arithmetic-jet}.
That composition is a typed map on the displayed local domain: physical
variables $(z,\sigma,t)$ remain distinct from the arithmetic variable $s$,
and the cutoff residual is retained before applying any Mellin functional.
The highest-diagonal identities therefore supply exact input to the existing
residue/regular-value hierarchy, but they do not imply a negative Weil
quadratic form, an off-critical zero, or a failure of the classical
Navier--Stokes theorem.  The full source series, its existence theorem and
all lower profile orders remain separate dependencies.
